\documentclass{ctexart}

\usepackage{multirow}
\usepackage{titlesec} %自u格式的宏包
\usepackage{picinpar,graphicx}
\usepackage{zhnumber}
\usepackage{float}
\usepackage{caption,subcaption}
\usepackage{amsmath}
\usepackage{amssymb}
\usepackage{fancyhdr}%页眉页脚
\usepackage{geometry}
\usepackage{titlesec} %自定义多级标题格式的宏包
\usepackage{amsmath}

\geometry{a4paper,left=2cm,right=2cm,top=2cm,bottom=2cm} %页边距

\title{职业生涯规划}
\author{游寅龙}
\date{\today}

\counterwithin{figure}{subsection}
\counterwithin{table}{subsection}

\renewcommand{\thefigure}{ \arabic{subsection}-\arabic{figure}}
\renewcommand{\thetable}{ \arabic{subsection}-\arabic{table}}


\pagestyle{plain}


\begin{document}
	\pagestyle{empty}
	\begin{figure}[h]
		
		\includegraphics[width=0.45\textwidth]{校名}
	
	\end{figure}
	
	\vspace*{10em}
	
	\begin{center}
		{\Huge 职业生涯规划与就业指导}\\
	\end{center}
	\vspace{14em}  % 垂直间距
	% 信息
	\begin{center}
		{\Large  % 这里的字号也可以用别的方式修改
			\makebox[4em][s]{学院}:\underline{\makebox[15em][c]{电气工程学院}}\\
			\makebox[4em][s]{专业}:\underline{\makebox[15em][c]{$22$级检测与仪器技术}}\\
			\makebox[4em][s]{姓名}:\underline{\makebox[15em][c]{游寅龙}}\\
			\makebox[4em][s]{学号}:\underline{\makebox[15em][c]{202211030001}}\\
			\makebox[4em][s]{指导老师}:\underline{\makebox[15em][c]{冯艳春}}\\
			\makebox[4em][s]{提交日期}:\underline{\makebox[15em][c]{2024年5月18日}}\\
		}
	\end{center}
	
	\clearpage  % 这一页结束后，页码重新开始
	\newpage
	
	\centering{\tableofcontents}
	\raggedright
	\setlength{\parindent}{2em}
	\setcounter{page}{0}
	\newpage
	\pagestyle{plain}
	\titleformat{\section}[block]{\zihao{3}\bfseries\centering}{第\zhnum{section}章 }{1em}{}[]
	\titleformat{\subsection}[block]{\zihao{-3}\mdseries\raggedright}{\zhnum{subsection}、}{1em}{}[]
	\titleformat{\subsubsection}[block]{\zihao{4}\bfseries\raggedright}{\arabic{subsubsection}.}{1em}{}[]
	\titleformat{\paragraph}[block]{\songti\zihao{5}\raggedright}{[\arabic{paragraph}]}{1em}{}[]
	
	\section{绪论}
		\subsection{课程学习情况}
			本学期我们学习了职业生涯规划与就业指导\uppercase\expandafter{\romannumeral2}。在本学期的学习过程中，我们跟随着老师的教学，学会了很多有用的知识，收获了很多的东西。本学期主要学习了三个大章节的内容，第一章节是大学生就业政策分析，老师一共讲了四个部分，第一部分是就业政策及发展沿革，第二部分是国家促进就业重大项目，第三部分是国家政策发展导向，第四部分是劳动关系有关政策、法规和法律。第二章节是专业发展与就业分析，主要分为两小节内容，分别是本专业专业概况、就业动态和本专业毕业生重点行业就业案例分析。第三章节主要是职业选择，分为三个小节的内容，第一小节是毕业的主要去向，第二小节是毕业的去向选择，第三小节是择业心理调适。跟随着老师的一步一步讲解，我们学到了很多和职业生涯规划与就业指导相关的知识，对我们今后的职业生涯规划与就业有着强有力的指导作用。
		\subsection{学习收获}
			此门课程中，我的个人思考收获还是有很多的，在选择自身职业之前，首先要考虑国家的政策，也就是国家在五年计划中重点发展与重点支持的方向，要充分了解未来就业的发展趋势，跟随国家的脚步。而后了解自身与自身专业，在国家重点发展方向中划定一个范围，选择自身适合的就业岗位，我们需要对岗位加快了解，而后审核个人因素、教育因素、环境因素和家庭因素是否满足自己可以选择该岗位。直到我们了解足够充分，才知道这是不是适合自身的岗位。\par
			除了就业之外，有很多毕业生选择了考研，我认为盲目地考研是不可取的，要看你能否用你的知识与能力驾驭选择的岗位和专业方向，或者是想要进大公司，亦或者是拿高薪……我认为考研和就业虽是两道道路，却仍有相似之处，考研的话，你的专业知识更加专业，也就是更加精准的匹配了就业领域与就业方向，本质是也是一种对职业的选择。所以，考研不仅仅是考研，也是对未来职业的一种确定。\par
			还有一条道路，就是自主创业，在目前的大环境下，世界的需求量下降，经济发展放缓，对于自主创业的人来说，现在的市场形势并不好，这就需要创业者有更多的资源和资金，但收益还不能得到保证。这并不适合大多数人，其中也包括我。\par
			此外，职业生涯规划教育能够让学生树立客观的就业目标大学生在大学期间就应当积极树立正确的就业目标，就业目标在一定的程度上反映了大学生的价值观、人生观，大学生在就业目标的前提下做好自己的职业规划。大学生进入大学深造其目的是通过不断升化自己，以期将来可以找到一份社会地位高、报酬多、契合自己兴趣并且能够实现个人价值的工作。所以说大学生就要树立正确的就业观，高校开展职业生涯规划教育能够让学生清楚地认识自己，认识社会工作的性质，并及早的树立正确、客观的就业目标。有了客观、正确的就业目标，大学生才能根据设定的目标合理规划自己的大学生活，进而在毕业时找到一份满意的工作。\par
			大学生对于职业生涯规划不能只考虑片面，要有整体性原则。大学生职业生涯规划要贯穿这一生，既要结合当前状况，又要顾及今后职业发展的空间，要对未来具有一定的预知性。大学生要从个人长期发展的角度出发制订职业生涯规划，明确这一生的使命是什么，让一切的奋斗都要沿着这个目标轨道行驶。由大及小，一步步前进，日积月累，最终实现总目标。俗话说得好：“计划不如变化。”世界每分每秒都在发生着变化，有一部分变化是可以提前预见并事先加以控制的，但是绝大多数变化是难以预见的。随着时间的推移，职业生涯规划总会受到这样或那样不确定性因素的影响，这就需要及时评估和调整职业生涯规划，以适应周围环境的变化，使之更切合个体发展的需要。\par
	
	\newpage
	
	\section{专业分析}
		\subsection{专业背景}
			测控技术与仪器专业培养德、智、体全面发展，具备测控技术、测控仪器与系统方面的基础知识与应用能力，具有扎实的理论基础和良好的业务素质，能够在国民经济各部门从事测量与控制领域的科技开发、仪器与系统的设计制造、应用研究、运行管理和经营销售等方面的高级工程技术人才。
		\subsection{专业发展和就业优势}
			测控技术与仪器专业在发展与就业方面拥有优势主要体现在以下几个方面：
			\begin{enumerate}
				\item 测控技术与仪器专业是一个多学科交叉融合的专业，它涵盖了测量技术、控制技术、计算机技术、通信技术等多个领域的知识。这种跨学科的特性使得该专业的学生在知识结构上更加全面，能够适应不同领域的工作需求，具备更强的竞争力和适应性。
				\item 测控技术与仪器专业具有广泛的应用领域。随着科技的不断进步和工业的快速发展，测控技术与仪器在制造业、航空航天、电力、医疗、环保等众多行业中扮演着重要角色。这些行业对测控技术与仪器专业人才的需求量大，为该专业的学生提供了丰富的就业机会。
				\item 测控技术与仪器专业注重实践能力的培养。在专业学习过程中，学生将接触到各种先进的测量仪器和控制系统，通过实验操作、课程设计、毕业设计等环节，不断提升自己的动手能力和解决实际问题的能力。这种实践能力的培养有助于学生在就业后迅速适应工作环境，承担工作任务。
				\item 测控技术与仪器专业还紧跟时代潮流，不断融入新技术、新方法。随着物联网、大数据、人工智能等技术的快速发展，测控技术与仪器专业也在不断更新教学内容和教学方法，培养学生的创新能力和适应能力。这使得该专业的学生在就业市场上更具竞争力。
			\end{enumerate}
		\subsection{专业发展和就业优势}
			测控技术与仪器专业在发展与就业方面拥有优势主要体现在以下几个方面：
			\begin{enumerate}
				\item 专业知识的广泛性可能导致深度不足。测控技术与仪器专业涵盖了多个领域的知识，虽然这使得毕业生在就业市场上具有广泛的适应性，但也可能导致他们在某一特定领域内的专业知识深度不够。这可能会限制他们在某些高级职位或专业领域的竞争力。
				\item 实践教学环节可能相对薄弱。测控技术与仪器专业是一门实践性很强的学科，需要学生通过实验、实习等方式来加深对理论知识的理解。然而，一些学校可能由于资源、资金等方面的限制，无法为学生提供充足的实践机会，这可能会影响学生的实践能力和就业竞争力。
				\item 行业发展的快速变化也可能带来挑战。随着科技的不断进步和产业的快速发展，测控技术与仪器领域也在不断更新换代。毕业生需要不断学习和适应新技术、新方法，才能跟上行业的发展步伐。这要求毕业生具备较强的学习能力和适应能力，否则可能会面临就业困难或职业发展受限的问题。
				\item 就业市场的竞争也可能较为激烈。由于测控技术与仪器专业的广泛应用和热门程度，吸引了大量学生报考。这使得该专业的毕业生在就业市场上可能会面临较为激烈的竞争。他们需要通过不断提升自己的综合素质和专业技能，才能在竞争中脱颖而出。
			\end{enumerate}
		\subsection{毕业去向}
			毕业生可到电气测量、检测检验、IT、物联网、人工智能、计量、精密工程、测量与控制等多领域的企事业单位从事信息检测、信息传输、信息处理、故障诊断、自动控制、仪器设计开发、系统集成、运行管理等方面工作。\par
			近三年毕业生就业去向有国家电网、华为、中兴通讯、海康威视、美的、TCL、广州视源等知名企业。
	\newpage
	\section{个人分析}
		\subsection{个人情况}
			在校期间，我认真学习了测控技术与仪器的相关本科课程，掌握了反馈控制理论，误差理论与数据处理等专业课程知识，成绩中等偏上，长期位于检测专业前十名，有望获得保研名额。\par
			在课后，我多次参与学科竞赛，获得优良成绩，包括但不限于：
			\begin{enumerate}
				\item 十八届全国大学生智能汽车竞赛智能视觉组省二等
				\item 十四届蓝桥杯全国软件和信息技术专业人才大赛省二等
			\end{enumerate}\par
			因此积累了一定的项目经验和团队合作能力，拥有丰富的实践经验和理论转实践的经验，非常合适未来与工作单位的技术接洽。\par
			本人在课内学习和竞赛实践中逐渐熟练掌握以下相关专业软件：
			\begin{enumerate}
				\item 办公类：Word、Excel、PowerPoint、Acrobat DC、Xmind、Typora
				\item 剪辑类：Adobe Photoshop、Adobe Premiere、KeyShot
				\item 机械类：SOLIDWORKS、AutoCAD、Bambu Studio
				\item 硬件类：立创EDA、Altium Designer
				\item 嵌入式类：Keil uVision5、IAR、OpenMV IDE、STM32CubeMX
				\item 软件类：Visual Studio、Visual Studio Code
				\item 学术类：TeXstudio、Multisim、MATLAB
			\end{enumerate}\par
			这些软件涵盖了多个行业多种职业，使得我的就业面极宽，可以是文工类、机械类、硬件类、嵌入式类、软件类的大多数本科生招生岗位。\par
			我对于技术创新和智能制造领域充满热情，希望能够通过自己的努力，为推动工业自动化和智能化发展做出贡献。我重视挑战性和学习成长，希望能够在工作中不断提升自己的技术水平和解决问题的能力。
		\subsection{家庭情况}
			我来自上海，有上海户口。父母健在，且没有离异。父母收入适中，能负担得起一家三口的所有合理开销。\par
			因此我的就业压力不大，家庭能承担我继续深造的开销。\par
			父母对于我未来的规划并没有过多的要求，我可以根据自我条件，时刻调整对于未来的规划。
	\newpage
	\section{毕业去向}
		\subsection{读研深造}
			\subsubsection{保送研究生}
				我的成绩目前是稳在前十，大约是第七、第八名。根据燕山大学保研规则，我需要在大学三年级下学期期末考试结束后处于专业前七才能被报送研究生。因此我需要将我的成绩稳中有进才能获取保研名额。\par
				获取方式有两个，其一是提升专业课成绩，我在过往的考试中发挥不稳定，成绩忽高忽低，如果能将发挥变稳定，那就可以提升$0.2-0.3$分的绩点；其二是参加竞赛获得国奖，我在过往的竞赛中均获得省级二等奖的成绩，仅差一步就能进入国赛，获得国家一次获得绩点加分。\par
				根据$20$级学长们的保研去向，检测最后一名的学长去向了上海理工大学，且因为他的保研名额是最后算分时候才获得的，并没有提前与导师交流，因此上海理工大学只是他最后临时的去向。因此我并不担心保研去向的学校，我的选择应当会在上海大学、上海理工大学等江浙沪的学校中选择。
			\subsubsection{考取研究生}
				因为我个人英语成绩并不理想，且对于考试发挥始终不稳定，因此我想我在将来的很长一段时间内都不会考虑这个选项。
		\subsection{入职就业}
			因为家庭关系，我会偏向于在上海找工作。因此所有数据均采用的是上海方面的数据，收入水平较高，生活成本也较高。
			\subsubsection{文工类}
				在此的文工类是对于偏向于交流类的工作，如：销售、管理等。\par
				虽然我对于办公剪辑类软件掌握较为熟练，但是由于我性格原因，我对人际交流不是很擅长，且这类工作专业难度低，可替代性强，平均工资低等缺点，只能作为我就业方面的下下策。
			\subsubsection{机械类}
				机械类的工作大致分为两类。\par
				第一类：仿真设计工程师，对应流体仿真等机械结构工程师，薪酬较高，但对应专业要求较高，这并不是测控专业善长的，难以与各高校机械本科生竞争。\par
				第二类：助理工程师，负责建模绘制，整理BOM表，难度较低，但平均工资也较低，税后收入在$5000$RMB上下，向上发展受限，难以通过时间有效向上攀爬。
			\subsubsection{硬件类}
				硬件类的工作大致分为电源工程师、汽车硬件工程师、信号工程师。对于这三类我都略有基础，且薪酬均在$1$万到$1.5$万左右，是本科生就业薪酬的中上等水准了，基本符合我的个人要求。\par
				但硬件类的工作大致有三个缺点：
				\begin{enumerate}
					\item 出差多，加班多，工作偶然性较强，在特定时间极其忙碌但是过段时间会极其空闲。
					\item 跳槽较为困难，隔行如隔山，电源工程师无法跳槽成为信号工程师，因为没有相关知识。
					\item 因为焊接使用助焊剂和松香等，会有一定的职业病。
				\end{enumerate}
			\subsubsection{嵌入式类}
				我个人认为随着应用软件（APP）市场的逐渐饱和，嵌入式将会成为下一个行业风口。\par
				虽然嵌入式包含嵌入式硬件和嵌入式软件，但此处的嵌入式类单指嵌入式软件，但相比如下述所述的软件类更偏向于硬件，主要是编写驱动，调试电路板，由于我既具有硬件的基础，也具备一定的算法编写能力，熟练掌握C/C++编程语言，因此可以胜任该职位。\par
				对于该职位，我目前能拿到的薪酬约在$7,000-10,000$RMB之间，如果想要提升的话需要掌握RTOS和Linux相关的知识，这部分的学习将占用我大量的时间，因此只有当我决心做这一行，我才会逐步的去学习。\par
				当然这一行的确定和软件类的大致一致，约有以下两点：
				\begin{enumerate}
					\item 生活极其不规律，以熬夜为主，想拿到较高薪酬，会涉及到996等工作常态。
					\item 会有程序员的通病，35岁转管理或是离职（当然随着目前出生人口的逐渐减少，是否会想如今一般大量转岗还有待考察）
				\end{enumerate}
			\subsubsection{软件类}
				软件类指的是应用软件（APP）的开发程序员，我能拿到的薪酬约在$8000$左右，但是随着应用软件市场的逐渐饱和，其实对于开发类型的程序员需求已经逐渐趋于饱和。\par
				市场会逐渐转变为程序维护员，因此随着人才需求的减少，就业形势并不容乐观。\par
				和机械类一样，软件类会作为我的备选。
				
	\newpage
	
	\section{个人计划}
		\subsection{未来去向}
			我不喜欢被某一条路所约束，所以我会为自己选择多条发展的方向。目前我的规划是三条路：\par
			\begin{enumerate}
				\item 保送研究生：我会尽一切可能获得绩点，以此增加我保研的成功率。
				\item 成为硬件工程师：我会以智能车队中学到的硬件知识为起点，扩展我对硬件的了解。
				\item 成为嵌入式工程师：我会增强RTOS和Linux方面的学习。
			\end{enumerate}\par
			以上的排序是我目前心中的排序，但我可能会在之后的关键时期，根据个人的实际发展做出进一步的规划。
		\subsection{未来计划}
			未来的计划大致分为三个时期。
			\subsubsection{初期}
				此处的初期约是指大二和大三学年、也就是近一年。\par
				在大二学年，我将参加完蓝桥杯和智能汽车竞赛。我计划在本学年内，这两个一类竞赛，我能进入国赛获得绩点加分。\par
				在大三学年，我会主攻学习，争取在学习上获得更好的成绩，以及通过大学英语六级考试，拿到合格证书。\par
				通过绩点加竞赛的方式力求获取保研名额。
			\subsubsection{中期}
				如果我有幸获得了保研名额，那我的中期计划会围绕研究生的学习来规划，争取在四年的研究生学习生涯中顺利毕业，找到一份优秀的工作。\par
				如果我不幸没有获得能保研的机会，那我的中期计划会围绕第一份工作规划，我倾向于靠近家附近的公司工作。我家旁边是张江高新技术园区，有一定的优秀创业公司，我的个人能力完全能胜任综合性能较强的研发工作。以此来提高初期的薪酬。
			\subsubsection{后期}
				后期大致是五六年后，也就是2030年前后，我已经快30岁，我会均衡生活与工作，可能会找一些略微轻松的工作，比如硬件工程师之类的。将一定的时间付诸于家庭。
		
\end{document}








